% BHU PYQ -- Manually transcribed & error-corrected
% Paper: BCF-222 : Financial Analysis
% Exam: B.Com. (Hons.) FMM Semester IV Examination, 2023-24

\documentclass[11pt,a4paper]{article}
\usepackage[a4paper,top=2.5cm,bottom=2.5cm,left=2.8cm,right=2.8cm,headheight=14pt]{geometry}
\usepackage{amsmath,amssymb}
\usepackage[T1]{fontenc}
\usepackage[utf8]{inputenc}
\usepackage{lmodern,microtype,fancyhdr,enumitem,hyperref,lastpage,parskip,booktabs}

\hypersetup{
    colorlinks=true,
    urlcolor=black,
    linkcolor=black,
    pdftitle={BCF-222: Financial Analysis -- BHU B.Com. (Hons.) FMM Sem IV 2023-24}
}

\pagestyle{fancy}
\fancyhf{}
\renewcommand{\headrulewidth}{0.4pt}
\renewcommand{\footrulewidth}{0.4pt}
\fancyhead[L]{\small   \textit{Banaras Hindu University}}
\fancyhead[R]{\small   \textit{BCF-222: Financial Analysis}}
\fancyfoot[C]{\small Page  \thepage\ of \pageref{LastPage}}


\newlist{parts}{enumerate}{2}
\setlist[parts,1]{label=(\roman*),leftmargin=*,itemsep=5pt,topsep=3pt}
\setlist[parts,2]{label=(\alph*),leftmargin=*,itemsep=3pt,topsep=2pt}

\newcommand{\pts}[1]{\hfill   \textit{[#1]}}


\begin{document}


\begin{center}
  {\large\bfseries BANARAS HINDU UNIVERSITY}\\[2pt]
  {
\normalsize B.Com. (Hons.) FMM --- Semester IV Examination, 2023-24}\\[6pt]
  \rule{0.8\linewidth}{0.5pt}\\[6pt]
  {\large\bfseries Paper No. BCF-222: Financial Analysis}\\[4pt]
  {
\normalsize Time: 3 Hours\hfill Maximum Marks: 70}
\end{center}

\rule{\linewidth}{0.4pt}

\smallskip



\noindent   \textit{Instructions: Answer all the questions. All questions are of equal marks (14 marks each).}

\bigskip




\noindent   \textbf{Question 1.}
What is meant by interpretation of financial statements? Explain its various components and discuss its importance.\pts{14}

\centerline{   \textbf{OR}}

Explain the method of preparing a comparative balance sheet and draw such a balance sheet on the basis of imaginary data.\pts{14}

\medskip\rule{\linewidth}{0.2pt}




\noindent   \textbf{Question 2.}
The following data are available from financial record of the company for the year 2022-23. On the basis of data, calculate the following ratios:

\begin{parts}
  \item Current Ratio
  \item Quick Ratio
  \item Inventory Turnover
  \item Payable Turnover
  \item Average Payment Period
\end{parts}\pts{14}


\begin{center}

\begin{tabular}{lr}
   oprule
   \textbf{Particulars} &   \textbf{Amount (Rs.)} \\
  \midrule
  Sales & 26,50,000 \\
  Sundry Creditors & 6,00,000 \\
  Bills Payable & 1,50,000 \\
  Purchase & 17,00,000 \\
  Purchase Returns & 37,500 \\
  Sundry Debtors & 6,00,000 \\
  Inventory & 4,00,000 \\
  Cash in Hand & 2,50,000 \\
  Cash at Bank & 10,00,000 \\
  ottomrule
\end{tabular}
\end{center}

\centerline{   \textbf{OR}}


\begin{parts}
  \item Capital investment in the company is Rs. 50,00,000 and its total turnover is 4 times of the capital investment. The gross profit margin on sales of the company is 5\%. What is the rate of return on investment of the company?\pts{7}
  \item What is debt equity ratio? Is there any standard debt equity ratio? Why is it necessary to maintain a balance between debt and equity invested in an enterprise?\pts{7}
\end{parts}

\medskip\rule{\linewidth}{0.2pt}


\newpage



\noindent   \textbf{Question 3.}
Golu Ltd. enterprises balance sheet on 31st March, 2021 and 2022 is as following:


\begin{center}

\begin{tabular}{lrr@{\qquad}lrr}
   oprule
   \textbf{Liabilities} &   \textbf{2021 (Rs.)} &   \textbf{2022 (Rs.)} &   \textbf{Assets} &   \textbf{2021 (Rs.)} &   \textbf{2022 (Rs.)} \\
  \midrule
  Share Capital & 40,00,000 & 52,00,000 & Goodwill & -- & 4,00,000 \\
  Reserve & 10,00,000 & 10,00,000 & Plant \& Machinery & 22,59,000 & 23,24,000 \\
  Profit \& Loss A/c & 7,93,800 & 8,24,400 & Building & 29,70,000 & 28,85,000 \\
  Sundry Creditors & 7,90,000 & 8,22,700 & Stock & 22,20,800 & 19,47,400 \\
  Bills Payable & 6,75,600 & 2,30,500 & Cash in Hand & 46,300 & 14,700 \\
  Bank Overdraft & 11,90,200 & -- & Sundry Debtors & 17,03,500 & 14,52,500 \\
  Provision for Taxation & 8,00,000 & 10,00,000 & Cash at Bank & 50,000 & 54,000 \\
  \midrule
   \textbf{Total} &   \textbf{92,49,600} &   \textbf{90,77,600} &   \textbf{Total} &   \textbf{92,49,600} &   \textbf{90,77,600} \\
  ottomrule
\end{tabular}
\end{center}

   \textbf{Additional Information:}

\begin{parts}
  \item An interim dividend of Rs. 5,20,000 was paid during 2022.
  \item The assets of another company were purchased for Rs. 12,00,000 payable in fully paid shares of the company. These assets include stock Rs. 4,32,800, machinery Rs. 3,67,200 and goodwill Rs. 4,00,000. In addition, miscellaneous items amounting to Rs. 1,13,000 were also purchased in connection with the plant purchase.
  \item An amount of Rs. 5,00,000 was paid during 2022 as Income Tax.
  \item The net profit (before tax) for the year 2022 was Rs. 12,50,000.
\end{parts}
Prepare a schedule of changes of working capital and the fund flow statement for the year 2022.\pts{14}

\centerline{   \textbf{OR}}

What do you understand by schedule of change in working capital? How does this schedule help in preparing the fund flow statement?\pts{14}

\medskip\rule{\linewidth}{0.2pt}




\noindent   \textbf{Question 4.}
What is the purpose of preparing a cash flow statement? How is it prepared? Explain with imaginary figures.\pts{14}

\centerline{   \textbf{OR}}

Explain the method of calculating `cash from operations' and clarify the difference between cash budget and cash flow statement.\pts{14}

\medskip\rule{\linewidth}{0.2pt}




\noindent   \textbf{Question 5.}
Write short notes on:

\begin{parts}
  \item Profit volume ratio and its various uses\pts{3.5}
  \item Margin of safety\pts{3.5}
  \item Fixed and variable cost\pts{3.5}
  \item Contribution\pts{3.5}
\end{parts}

\centerline{   \textbf{OR}}

The following is the cost information relating to manufacturing of a product:

\begin{center}

\begin{tabular}{ll}
  Selling Price & Rs. 20 per unit \\
  Trade Discount & 5\% of selling price \\
  Material Cost & Rs. 6 per unit \\
  Labour Cost & Rs. 4 per unit \\
  Overheads: & \\
  \quad Fixed Cost & Rs. 20,000 \\
  \quad Variable Cost & 100\% of labour cost \\
\end{tabular}
\end{center}
You are required to calculate the following:

\begin{parts}
  \item Break Even Point\pts{7}
  \item Profit if sales are 10\% and 15\% above the break even point.\pts{7}
\end{parts}

\vfill
\rule{\linewidth}{0.4pt}

\begin{center}\small
\href{https://bhu-exam.vercel.app/}{Click here to visit our website}\\[4pt]
\href{https://me-aryan.vercel.app/}{Click here to visit our contributor}\\[4pt]
\href{https://quashan-me.vercel.app/}{Click here to visit our contributor}
\end{center}

\end{document}
